\section{From a seed energy gain to packet estimates}
\label{sec:packet-estimates}
We now supply the passage from local energy and chart estimates to the
packet price used in Section~\ref{sec:mixed-activity-prices}. There are
three distinct costs. Extending the energy from a seed core to its full
support introduces a small annular error. The density and normalization
factors introduce a polynomial logarithmic expense. Differentiating the
packet introduces a further polynomial factor. All three are paid from
the same energy gain per seed.

Keep the scales $h,P,d,\tau,B$ of
\eqref{eq:warning-scales}, and fix the block length, chart, shell width
and all other local parameters before sending $x$ to infinity. All
constants below are independent of the lattice period. For clarity,
say that $r(x)$ has upper exponent $a$ if, for every $b>a$,
$|r(x)|\le h(x)^b$ eventually. This permits subpower losses. Products
add upper exponents, and a fixed finite sum preserves a common upper
exponent: in each case choose a strictly smaller exponent than the
one finally required and absorb fixed constants using $h\to\infty$.
In particular $B$, $d$ and every fixed real power of $P$ have upper
exponent zero, while $\tau$ has upper exponent $-1/2$.

\subsection{Extending a core energy bound}
Consider a packet with $n$ seeds. In each signed channel let
$D_2\subseteq D_{200}$ be the sets of local energy addresses within
$2B$ and $200B$ of its seeds. Their sizes satisfy
\[
 |D_{200}|\le n\,b_{200B+1}\,s_L,
\]
where $s_L$ bounds the number of local energy addresses per coarse
site. Let $e_p$ be the assigned energy at address $p$ in that channel.
Suppose the core energy is $\sum_{p\in D_2}e_p$ and the annular terms
satisfy $e_p\ge-r(x)$ on $D_{200}\setminus D_2$, with $r(x)\ge0$.
Then
\[
 \sum_{p\in D_{200}}e_p
 \ge\sum_{p\in D_2}e_p-r(x)b_{200B+1}s_L n.
\]
Average this inequality over signed channels with their normalized
weights. If the averaged core energy is at least $\mu\tau^2n$, the
averaged extended energy $Q_A$ obeys
\begin{equation}\label{eq:packet-assigned-energy}
 Q_A\ge(\mu\tau^2-E(x))n,
 \qquad E(x)=r(x)b_{200B+1}s_L.
\end{equation}
No energy gain from one channel is multiplied by a factor from another:
the average is taken on the common configuration and common packet.

The local inputs in this application are a positive seed charge
$\mu$, the displayed core bound, and an annular remainder with upper
exponent $-3/2$. The latter is proved in
Section~\ref{sec:packet-annulus} by comparing the fine and predictor
energies with their extracted affine squares. On an address where no quadratic term
was extracted, the fine and predictor energies are nonnegative;
only the one-sided lower bound is needed there. On selected addresses,
the local Taylor estimates bound the difference from the extracted
quadratic energy in absolute value. The protected $304B$ fluctuation
region and the $300B$ mask conditions allow those local estimates to
be used throughout the $200B\setminus2B$ annulus. Multiplying $r$ by the
fixed-density ball count preserves its upper exponent, so $E$ also
has upper exponent $-3/2$.

\subsection{The packet density and its normalization}
Let $A$ denote the packet. Its logarithmic density has the form
\begin{equation}\label{eq:packet-log-density}
 \Phi_A=-\beta(Q_A+Q_A^{\rm mean})
       +\sum_{i\in V_{200}(A)}c_i
       +\sum_{i\in V_{304}(A)}j_i
       -N_A\log m_\beta.
\end{equation}
Here $Q_A^{\rm mean}\ge0$ is the extracted mean-square contribution;
$c_i$ are the local Gaussian normalization corrections; $j_i$ are the
protected chart corrections; and $N_A$ counts the free chart bonds
within distance $10B$ of the packet. The bond sets and this count
satisfy, with constants fixed at $L$,
\[
 |V_{200}(A)|\le C_{200}d^4n,\qquad
 |V_{304}(A)|\le C_{304}d^4n,\qquad
 N_A\le C_{10}d^4n.
\]
Cover each seed neighbourhood by a lattice ball. There are
$L^4\#\{\text{axes}\}$ fine bonds per coarse site, so summing these
bounds over seeds gives the claimed counts. Overlaps reduce the count.

The one-bond Gaussian normalization is
\[
 m_\beta=c_{\rm chart}(\sqrt{2\pi}\,\beta^{-1/2})^q,
 \qquad c_{\rm chart}>0,
\]
where $q$ is the colour dimension. Thus
\[
 -\log m_\beta
 =-\log c_{\rm chart}-\frac q2\log(2\pi)
                        +\frac q2\log\beta
 \le C_{\rm norm}d
\]
for $0<\beta\le2h$ and sufficiently large $x$. The bound holds even
when $\beta$ is very small, since only an upper bound on $\log\beta$
is used. Suppose also $|c_i|\le C_{\rm cov}$, $|j_i|\le1$, and that
the likelihood-ticket product is bounded by
$\exp(C_{\rm tag}d^4n)$. These are uniform local analytic bounds.
The other packet masks, enforcement factors and contractions take
values in $[0,1]$. Dropping the nonpositive mean-square term in
\eqref{eq:packet-log-density}, and using $d\ge1$, gives
\begin{equation}\label{eq:packet-normalization}
 \Phi_A+C_{\rm tag}d^4n
 \le-\beta Q_A+K d^5 n
\end{equation}
for a fixed $K\ge0$.

Here the ticket bound follows directly from the normalized tag
probabilities. Let $N_f$ and $N_p$ be the numbers of frame and
predictor tags in the fixed patch. A site likelihood is $N_f$ times
its tag probability, and a bond likelihood is $N_p$ times its tag
probability. The probabilities are nonnegative and at most one:
in the minimizing frame branch they are uniform on the nonempty
set of minimizers, in the fallback branch they are uniform on all
tags, and the predictor weights are nonnegative with sum one.
Thus each likelihood is at most $N_f\le e^{N_f}$ or
$N_p\le e^{N_p}$, respectively. The radius-$r$ seed neighbourhood
contains at most $n b_{r+1}$ sites, and the bonds incident to it
number at most $n b_{r+2}\#\{\mathrm{axes}\}$. Multiplying one
factor at each such physical site and bond gives
\[
 \mathrm{tickets}_r\le
 \exp\!\left[n\{b_{r+1}N_f+
 b_{r+2}\#\{\mathrm{axes}\}N_p\}\right].
\]
For $r=300B$ this is the asserted $\exp(C_{\rm tag}d^4n)$
bound. A tag used by several local tests is still charged only once
in this product. There is no independent tag resampling at each use.

A packet is set to zero unless its seed support is contained in the
sector seed support and its own flag set agrees with the actual coarse
flags throughout its $300B$ neighbourhood. On a healthy packet
neighbourhood, where there are no actual coarse flags, this flag test is
equivalent to the packet's flag component being empty. The latter is
the eligibility rule of the locally smooth stock extension used in the local
jet argument; its equality with the physical packet also uses the
funded room and manifest-sector conditions and a healthy neighbourhood
through radius at least $500B$. Indeed, every packet flag lies in its
own support and hence in its $300B$ neighbourhood, so the flag
agreement forces the packet flag set to be empty when that neighbourhood
is healthy; the converse follows because both flag sets then vanish
there. For the remaining factors, the $500B$ health region contains the
$200B$ energy addresses together with their clean read neighbourhoods,
as well as the $300B$ and $304B$ mask regions. The clean-hull agreement
of the clipped and unmasked rows and the manifest-sector mean identity
therefore identify the two exponents; the held fine field, likelihoods,
masks and normalization factors are the same in both products.
Neither rule requires the sector
store's flag component to be globally empty. When the packet value is
nonzero, every retained mask is nonzero. The mask
support conditions then give the local hypotheses of
\eqref{eq:packet-assigned-energy}, including the control of the hot
fluctuation bonds. Therefore
\[
 |\mathcal P_A|
 \le\exp\{\Phi_A+C_{\rm tag}d^4n\}
 \le\exp\{((K d^5+\beta E)-\beta\mu\tau^2)n\}.
\]
Since $\beta\le2h$ and $E$ has upper exponent $-3/2$, the product
$\beta E$ tends to zero and is eventually at most $1\le d^5$.
We have proved
\begin{equation}\label{eq:packet-value-price}
 |\mathcal P_A|\le
 \exp\{(C d^5-\beta\mu\tau^2)n\},\qquad C=K+1.
\end{equation}
The estimate also holds for packets which vanish, since its right-hand
side is positive. This argument retains the Gaussian normalization
$m_\beta$; it does not replace the original density by an unnormalized
free density.

\subsection{Derivatives at the edge of a mask}
For a $C^j$ map $F$ defined near a point $y$ in a finite-dimensional
real normed space, write
\[
 F\in\mathcal B_j(U,v;y)
 \quad\hbox{if}\quad
 \|D^kF(y)\|\le Uv^k\quad(0\le k\le j),
 \qquad U,v\ge0.
\]
The zero-order bound is included. The product rule implies
\begin{equation}\label{eq:packet-derivative-product}
 F\in\mathcal B_j(U,v;y),\quad G\in\mathcal B_j(V,w;y)
 \ \Longrightarrow\ FG\in\mathcal B_j(UV,v+w;y).
\end{equation}
Indeed the norm of the $k$th derivative is bounded by
$UV\sum_{r=0}^k\binom kr v^rw^{k-r}=UV(v+w)^k$.
Repeated application gives a product bound with multiplied amplitudes
and added derivative scales.

For the packet's stock-extension variation, the underlying fine
configuration in its mask and likelihood factors is held fixed; the coarse field and its
compensating fluctuation vary coherently. This matters at the boundary
of a cutoff. If any of the retained mask products is identically zero
on a neighbourhood of $y$, the extension and all its derivatives through
order $j$ vanish at $y$. In the complementary case, every relevant
individual mask has points of nonzero value arbitrarily close to $y$.
Continuity of its argument therefore places that argument in the
\emph{closed} support of the mask at $y$. For example, a cutoff which
vanishes when $\|u\|>2a\tau$ implies $\|u(y)\|\le2a\tau$ in this
case. The same argument gives the closed fine and predictor threshold
inequalities. It does not claim that the mask value at $y$ is nonzero.
The scalar estimate \eqref{eq:packet-normalization} and the core
estimate must hold on these closed conditions; this is the closed
support version of the local input used below.

For the local jet estimate, use the locally smooth stock extension described
above; its seed-eligibility indicator is fixed as the coarse field
varies. Separate its expression as $e^{\phi(v)}H(v)$, apart from that
indicator. The local differentiated estimates give
\begin{equation}\label{eq:packet-local-jet-inputs}
 \begin{split}
 \phi-\phi(y)&\in\mathcal B_j(b(x)n,1;y),\\
 H&\in\mathcal B_j(e^{\ell n},c(x)n;y),\\
 e^{\phi(y)+\ell n}&\le e^{(C d^5-g)n},
       \qquad g=\beta\mu\tau^2,
 \end{split}
\end{equation}
where $b,c\ge0$, $b$ has upper exponent $1$ and $c$ upper exponent
zero. Here $\ell$ is the ticket expense per seed. The first bound is
\eqref{eq:action-packet-log-jets}. To check the bound for $H$, first
consider a coarse noise cutoff. Its argument is the squared chord
of $V_e Q(U)_e^{-1}$ divided by $(W\tau)^2$, with $W>0$ fixed
and $U$ held fixed. If the chord is strictly larger than $2W\tau$,
the inner cutoff is locally zero. Otherwise the normalized matrix
displacement has bounded value and derivatives of upper exponent
zero: its positive derivatives are bounded using $v/\tau$, where
$v$ has upper exponent $-1/2$ and $\tau^{-1}$ has upper exponent
$1/2$. Composition with the fixed smooth cutoff on $[0,4]$ gives
derivative scale of upper exponent zero. Its value is in $[0,1]$.
The complementary cutoff $1-\chi$ has the same property: its value
is still at most one, and its positive derivatives differ only by
a sign. At a boundary point use the closed chord inequality, as
above. A fixed enlargement of the derivative scale absorbs the
finite-order composition constants while keeping amplitude one.

The protected bond cutoffs have the corresponding bound by
Section~\ref{sec:haar-chart} and the compensated increment estimate.
Fine masks, contractions and likelihoods evaluated at the held fine
configuration are constants. The first two have absolute value at
most one; the product of likelihoods has the stated bound
$e^{\ell n}$. Apply \eqref{eq:packet-derivative-product} to all these
factors. There are at most $n b_{300B+2}\#\{\mathrm{axes}\}$ coarse
noise cutoffs and $n b_{304B+2}\rho_L$ protected bond cutoffs,
where $\rho_L$ is the fine-bond density. Their derivative scales
add, giving $c(x)n$ with upper exponent zero, as asserted.
The last inequality in
\eqref{eq:packet-local-jet-inputs} is the closed support version just
explained. Thus the differentiated local action, covariance and chart
bounds proved in the preceding sections supply the required estimates.

We give the remaining derivative calculation. The repeated chain rule
for $e^\phi$ has one term for each partition $\pi$ of the $k$ derivative
arguments:
\[
 D^k(e^\phi)(y)[v_1,\ldots,v_k]
 =e^{\phi(y)}\sum_\pi\prod_{C\in\pi}
 D^{|C|}\phi(y)[v_i:i\in C].
\]
This identity follows by induction: differentiation either adds the
new argument as a singleton block or inserts it into one existing
block. Each partition appears exactly once. There are at most $k!$
partitions: send each partition to the permutation whose cycles list
each block in increasing order. This map is injective. For unit
arguments, \eqref{eq:packet-local-jet-inputs} therefore gives, for
$1\le k\le j$,
\[
 \|D^k(e^\phi)(y)\|
 \le e^{\phi(y)}k!(1+bn)^k
 \le e^{\phi(y)}\{J_j(1+bn)\}^k,
 \qquad J_j=\max\{1,j!\}.
\]
The zero-order bound is equality. Using
\eqref{eq:packet-derivative-product} and $n\ge1$, the packet belongs to
\[
 \mathcal B_j\bigl(e^{(C d^5-g)n},a(x)n;y\bigr),
 \qquad a(x)=J_j(1+b(x))+c(x).
\]
The function $a$ has upper exponent $1$, so $a(x)\le h^2$ eventually.
As $d\ge1$ and $d\ge\log h$ eventually, for $n\ge1$,
\[
 a(x)\le e^{2d^5n},\qquad n\le e^{d^5n},\qquad
 (a(x)n)^j\le e^{3j d^5n}.
\]
Enlarge $a$ to at least one if necessary. Every derivative of order
at most $j$ is now bounded by the same number:
\begin{equation}\label{eq:packet-jet-price}
 \|D^k\mathcal P_A(y)\|
 \le\exp\{((C+3j)d^5-g)n\},\qquad 0\le k\le j.
\end{equation}
Thus finite differentiation changes the fixed coefficient of the
logarithmic expense; it does not consume a fixed fraction of the
energy gain. This argument applies separately at every point of the
allowed variation, with the same constants.

\subsection{Absorbing the expense and the packet count}
Here is the quantitative comparison used to turn these estimates into
small connected-support prices. Suppose $\mu>0$ and
$\beta\ge u h$ for some fixed $u>0$. Define
\[
 M=\mu u P^2,\qquad q_*=e^{-M/2}.
\]
Since $h\tau^2=P^2$, the gain satisfies $g=\beta\mu\tau^2\ge M$.
For any fixed $D\ge0$, integer $r\ge0$ and $z<2$, the monitor
comparison \eqref{eq:warning-monitor-pays} gives
\begin{equation}\label{eq:packet-expense-absorption}
 Dd^rP^z\le M/16
 \quad\hbox{eventually}.
\end{equation}
Indeed apply that comparison with exponent $2-z$ and coefficient
$16D/(\mu u)$, and multiply by $\mu uP^z/16$.

Fix a packet support radius parameter $\ell$ and a support exponent
$\alpha\ge0$. Let $H_B$ and $D_B$ be the root and adjacency counts
from Section~\ref{sec:mixed-activity-prices}, and let $c_\ell$ be its
cover cost per seed. All are polynomial in $B$ except the fixed
$c_\ell$. Consequently, for sufficiently large $x$,
\[
 H_B\le e^{M/16},\qquad D_B^2\le e^{M/16},\qquad
 \alpha c_\ell\le M/16,\qquad 2\le M/16.
\]
For example, bound both counts by fixed powers of $d$, use
\eqref{eq:packet-expense-absorption}, and then $t\le e^t$.
Consider a nonnegative packet price
$C_0e^{(s-g)n}$ satisfying
\[
 s\le Dd^rP^z,\qquad 0\le C_0\le e^{Dd^rP^z}.
\]
Then $C_0\le e^{M/16}$ and the counting ratio obeys
\[
 D_B^2e^{s-g+\alpha c_\ell}
 \le e^{M/16+M/16-M+M/16}\le q_*.
\]
In particular $q_*\le1/2$. The packet counting estimate already proved
therefore bounds the site sum over packets with at least $n_0$ seeds by
\[
 H_BC_0\frac{q_*^{\max(1,n_0)}}{1-q_*}
 \le\exp\{-M\max(1,n_0)/4\}.
\]
To verify the last inequality, put $r_0=\max(1,n_0)$. Since
$(1-q_*)^{-1}\le2\le e^{M/16}$, its logarithm is bounded above by
$3M/16-Mr_0/2\le-Mr_0/4$.
This applies to \eqref{eq:packet-value-price} and
\eqref{eq:packet-jet-price} with $z=0$ and $r=5$. Hence once the
stated local energy, chart and derivative inputs have been established,
their full packet and support sums are small uniformly in the period.
The scalar comparison selects a threshold after all fixed constants;
it does not by itself construct the schedule that reaches that
threshold at each renormalization step.
